\documentclass[11pt,a4paper]{article}

% ── Packages ──────────────────────────────────────────────────────
\usepackage[margin=2.2cm]{geometry}
\usepackage{fontspec}
\setmainfont{DejaVu Serif}
\setsansfont{DejaVu Sans}
\setmonofont{DejaVu Sans Mono}
\usepackage{xcolor}
\usepackage{enumitem}
\usepackage{hyperref}
\usepackage{titlesec}
\usepackage{booktabs}
\usepackage{longtable}
\usepackage{colortbl}
\usepackage{tabularx}
\usepackage{amsmath,amssymb}
\usepackage{microtype}

% ── Colours ───────────────────────────────────────────────────────
\definecolor{highlightblue}{RGB}{0,51,102}
\definecolor{softblue}{RGB}{220,230,241}
\definecolor{darkgreen}{RGB}{0,100,0}
\definecolor{softgreen}{RGB}{230,245,230}
\definecolor{softred}{RGB}{245,220,220}

% ── Section formatting ────────────────────────────────────────────
\titleformat{\section}{\Large\bfseries\color{highlightblue}}{}{0em}{}[\vspace{-2pt}\rule{\textwidth}{0.4pt}]
\titleformat{\subsection}{\large\bfseries\color{highlightblue!80!black}}{}{0em}{}
\titleformat{\subsubsection}{\normalsize\bfseries\color{highlightblue!60!black}}{}{0em}{}
\pagenumbering{gobble}

% ── Custom commands ───────────────────────────────────────────────
\newcommand{\goodenough}[1]{\textcolor{darkgreen}{\textit{#1}}}
\newcommand{\toomuch}[1]{\textcolor{red}{\textit{#1}}}
\newcommand{\done}{\textcolor{darkgreen}{$\checkmark$}}
\newcommand{\todo}{\textcolor{orange}{$\circ$}}

\newcolumntype{L}[1]{>{\raggedright\arraybackslash}p{#1}}

% ── Document ──────────────────────────────────────────────────────
\begin{document}

\title{\vspace{-1cm}\textcolor{highlightblue}{\textbf{opto-sim Research Roadmap}}}
\author{Physical-Layer Fiber-Optic Simulator}
\date{}
\maketitle

\noindent A tiered publication plan with effort estimates,
experimental-validation requirements, and scope boundaries.

\vspace{0.5cm}

% ══════════════════════════════════════════════════════════════════
\section{Perspective}

opto-sim is a \textbf{general-purpose, open-source, validated physical-layer
fiber-optic simulator}. The complex-envelope electric field
($\overline{|\mathcal{E}|^2} = \text{Watts}$) is the single source of truth.
BB84 QKD is one downstream application protocol, not the defining purpose.

This roadmap is organised around \textbf{four papers}. Tier~1 validates the
simulator itself against published experimental data. Tiers~2--4 are QKD-focused
and depend on Tier~1 being published.

% ══════════════════════════════════════════════════════════════════
\section{Publication Plan}

\vspace{0.3cm}
{\small
\begin{tabularx}{\textwidth}{c L{3.5cm} L{2.8cm} c L{2.5cm}}
\toprule
\# & Paper & Target & When & Depends on \\
\midrule
A & Software validation: validated physical-layer simulator & OE\,/\,JLT\,/\,SoftwareX & Now & --- \\
B & QKD environmental sensitivity: temperature, bend, channel & OE\,/\,JLT\,/\,EPJ~QT & After A & A \\
C & Physical-layer security: side-channel analysis & IEEE~T-IFS\,/\,NPJ~QI\,/\,PRX~Q. & After B & A+B \\
D & Full QKD system: decoy-state, SKR, release & Nature Comms\,/\,PRX~Q. & Optional & A+B+C \\
\bottomrule
\end{tabularx}}

% ══════════════════════════════════════════════════════════════════
\section{Tier 0 --- Infrastructure (COMPLETE)}

\begin{tabularx}{\textwidth}{l L{6cm} l}
Effort: 1--2 weeks & Deliverable: pytest suite + seeded reproducibility + 48 tests & Paper: None (scaffolding) \\
\end{tabularx}

\noindent\textbf{Status:} All done. No further work planned.

% ══════════════════════════════════════════════════════════════════
\section{Tier 1 --- Software Validation Paper}

\textbf{Goal:} Publish a validated physical-layer fiber-optic simulator that
reproduces published experimental data across all components.

\textbf{Novelty:} No open-source simulator validates CD, PMD, birefringence,
attenuation, laser RIN, phase noise, MZM transfer, and APD response against
\textit{both} analytic theory and published experimental/datasheet curves.

% ── 1.1 ───────────────────────────────────────────────────────────
\subsection{1.1  Fiber Channel (DONE)}

\begin{tabularx}{\textwidth}{L{3.5cm} L{4cm} c L{3cm}}
\toprule
Validation & Reference & Error & Script \\
\midrule
CD: Gaussian pulse broadening & Agrawal Fig.~2.6 & 0.0000\% & \texttt{validate\_cd.py} \\
PMD: DGD histogram vs Maxwellian & Razavi Fig.~2.11 & KS $p=0.31$ & \texttt{validate\_pmd.py} \\
Attenuation: $\alpha=0.182$\,dB/km @ 1550\,nm & SMF-28 spec & 0.0000\% & \texttt{validate\_attenuation.py} \\
Birefringence: $\Delta n \propto (r_f/R)^2$ & Yuan Fig.~1 & 0.0000\% & \texttt{validate\_birefringence.py} \\
\bottomrule
\end{tabularx}

These prove the implementation matches the underlying analytic formula to
machine precision. They are \textit{necessary} but not \textit{sufficient}
for a software paper---they verify coding correctness, not physical fidelity
against real-world data.

% ── 1.2 ───────────────────────────────────────────────────────────
\subsection{1.2  Fiber --- Experimental Overlay (NEW)}

\subsubsection{D($\lambda$) dispersion curve --- Corning SMF-28 datasheet}

\begin{tabularx}{\textwidth}{L{3cm} L{10cm}}
\toprule
What to add & Sellmeier coefficients for fused silica (Malitson~1965)
$\rightarrow n(\lambda) \rightarrow \mathrm{d}^2n/\mathrm{d}\lambda^2
\rightarrow D_{\text{material}}(\lambda) + D_{\text{waveguide}}$ \\
Effort & $\sim$2 sessions \\
Good enough & $|D_{\text{sim}} - D_{\text{datasheet}}| < 1$\,ps/(nm$\cdot$km) across O--L bands \\
Too much & Full waveguide mode solver; temperature-dependent Sellmeier \\
Priority & \textbf{High} --- $D(\lambda)$ is the standard fiber datasheet curve \\
\bottomrule
\end{tabularx}

\subsubsection{$\alpha(\lambda)$ attenuation spectrum --- Corning SMF-28 datasheet}

\begin{tabularx}{\textwidth}{L{3cm} L{10cm}}
\toprule
What to add & Rayleigh ($C_R/\lambda^4$) + IR tail + UV tail + OH$^-$ peak at 1383\,nm \\
Effort & $\sim$1 session \\
Good enough & $|\alpha_{\text{sim}} - \alpha_{\text{datasheet}}| < 0.02$\,dB/km outside OH peak \\
Too much & Temperature-dependent $\alpha$; radiation-induced loss; bend-loss spectrum \\
Priority & \textbf{High} --- attenuation spectrum is the second standard fiber curve \\
\bottomrule
\end{tabularx}

% ── 1.3 ───────────────────────────────────────────────────────────
\subsection{1.3  Laser Validation (NEW)}

\subsubsection{RIN spectrum --- $S_{\text{RIN}}(f)$ vs relaxation-oscillation model}

\begin{tabularx}{\textwidth}{L{3cm} L{10cm}}
\toprule
What to add & Validation script: RIN time series $\rightarrow$ Welch PSD $\rightarrow$
Coldren Eq.~5.3.38 overlay \\
Effort & $\sim$1 session \\
Good enough & Peak frequency $\pm10\%$; roll-off $1/f^2$; DC level matches RIN$_0$ \\
Too much & Shot-noise-calibrated high-resolution RIN; $1/f$ noise \\
Priority & \textbf{High} --- relaxation-oscillation model is a unique simulator feature \\
\bottomrule
\end{tabularx}

\subsubsection{Phase noise --- Lorentzian linewidth verification}

\begin{tabularx}{\textwidth}{L{3cm} L{10cm}}
\toprule
What to add & Simulate field $\rightarrow$ complex-envelope PSD $\rightarrow$
fit Lorentzian, compare FWHM to input $\Delta\nu$ \\
Effort & $\sim$1 session \\
Good enough & Fitted FWHM within $\pm10\%$ of input linewidth \\
Too much & Gaussian line-shape decomposition; $1/f$ divergence \\
Priority & \textbf{High} --- linewidth validation is essential \\
\bottomrule
\end{tabularx}

% ── 1.4 ───────────────────────────────────────────────────────────
\subsection{1.4  MZM Validation (NEW)}

\begin{tabularx}{\textwidth}{L{4cm} L{3.5cm} L{3cm} L{2.5cm}}
\toprule
Validation & What to add & Effort & Good enough \\
\midrule
$V_\pi$ from crystal params & Report for default geometry & None (test exists) & $\pm5\%$ of literature \\
Extinction ratio & Plot $E_{\text{out}}$ vs $V$ & None (test exists) & Matches user ER \\
Push-pull vs single-drive chirp & Compare output phase & $\sim$1 session & Chirp $\propto$ 0.5$V_\pi$ \\
\bottomrule
\end{tabularx}

\textbf{Priority: Medium.} Well-tested model; one validation script to publish.

% ── 1.5 ───────────────────────────────────────────────────────────
\subsection{1.5  APD Validation (NEW)}

\begin{tabularx}{\textwidth}{L{4cm} L{3.5cm} L{3cm} L{2.5cm}}
\toprule
Validation & What to add & Effort & Good enough \\
\midrule
Responsivity $R(\lambda)$ & Spectral $\eta(\lambda)$ for InGaAs; plot vs datasheet & $\sim$1 session & Peak $\pm10\%$, shape matches \\
Shot-noise scaling & $I_{\text{noise}}^2 = 2eIB\cdot F$ & None (test exists) & Numerical tolerance \\
\bottomrule
\end{tabularx}

\textbf{Priority: Medium.}

% ── 1.6 ───────────────────────────────────────────────────────────
\subsection{1.6  System Demonstration}

\begin{tabularx}{\textwidth}{L{5cm} L{3cm} L{3cm}}
\toprule
Demonstration & Status & Effort \\
\midrule
BB84 QBER under CD+PMD vs distance & Existing graphic & 0 \\
Full link budget: laser $\rightarrow$ fiber $\rightarrow$ detector & New & $\sim$1 session \\
\bottomrule
\end{tabularx}

\textbf{Priority: Medium} --- completes the system story.

% ── 1.7 ───────────────────────────────────────────────────────────
\subsection{1.7  Tier 1 --- Scope Boundaries}

\subsubsection*{IN SCOPE}
\begin{itemize}[nosep]
\item Fiber: $D(\lambda)$ dispersion curve, $\alpha(\lambda)$ attenuation spectrum
  (Sellmeier + Rayleigh + OH$^-$)
\item Laser: RIN spectrum plot, Lorentzian linewidth verification
\item MZM: $V_\pi$, extinction, chirp-mode validation scripts
\item APD: Spectral responsivity, shot-noise validation
\item System: BB84 QBER demonstration, link budget
\end{itemize}

\subsubsection*{OUT OF SCOPE}
\begin{itemize}[nosep]
\item Nonlinear effects (Kerr, SPM, XPM, FWM) --- 3+ months, separate paper
\item EDFA / optical amplifier --- not needed for passive-fiber validation
\item PDL --- second-order effect, no reference data
\item Splice / connector loss --- negligible publication value
\item Temperature-dependent Sellmeier --- belongs in Tier~2
\item Waveguide mode solvers --- defeats the purpose of a macroscopic simulator
\item More protocols (E91, MDI-QKD, TF-QKD) --- strong diminishing returns
\item Quantum security analysis --- belongs in Tier~3
\end{itemize}

\subsubsection*{Effort Summary}

\begin{tabularx}{\textwidth}{L{3cm} L{2.5cm} L{3.5cm} L{2cm}}
\toprule
Sub-tier & Sessions & New code & Status \\
\midrule
1.1 Fiber (analytic) & --- & --- & \done \\
1.2 Fiber (experimental) & $\sim$3 & Sellmeier, $\alpha(\lambda)$ & \todo \\
1.3 Laser & $\sim$2 & RIN/linewidth scripts & \todo \\
1.4 MZM & $\sim$1 & Validation script & \todo \\
1.5 APD & $\sim$1 & Spectral $\eta(\lambda)$ model & \todo \\
1.6 System & $\sim$1 & Link budget script & \todo \\
\midrule
\textbf{Total} & \textbf{$\sim$8 sessions} & & \\
\bottomrule
\end{tabularx}

Estimated calendar time: \textbf{3--4 weeks} combined coding + paper writing.

% ══════════════════════════════════════════════════════════════════
\section{Tier 2 --- QKD Environmental Sensitivity (DEFERRED)}

\textbf{Requires:} Tier~1 published first.

\textbf{Goal:} Publish a journal paper on how temperature and bend radius affect
QKD system performance.

\subsection*{Tasks}
\begin{enumerate}[nosep]
\item Temperature sweep (0--50$\,^{\circ}$C, 25\,km spans) $\rightarrow$
  $\Delta n(T)$, beat length, QBER
\item Bend sensitivity ($R = 10$--100\,mm) $\rightarrow$ $\Delta n(R)$,
  coupled $T \times R$ regimes
\item Time-varying channel (1$\,^{\circ}$C/min ramp) $\rightarrow$ QBER time series
\item Aerial vs buried comparison (ASHRAE weather data) $\rightarrow$ QBER variance
\end{enumerate}

\subsection*{Scope Boundaries}

\begin{tabularx}{\textwidth}{L{7cm} L{7cm}}
\toprule
IN & OUT \\
\midrule
Temperature-dependent birefringence and phase drift & Temperature-dependent Sellmeier \\
Static and slow-ramp temperature profiles & Thermo-mechanical stress modeling \\
Single-span fiber (25\,km) & Multi-span with amplifiers \\
BB84 phase-encoding & Other protocols \\
Bend-induced birefringence & Bend-induced attenuation \\
\bottomrule
\end{tabularx}

\textbf{Effort: $\sim$6 weeks.} Stop if QBER delta $< 0.5\%$.

% ══════════════════════════════════════════════════════════════════
\section{Tier 3 --- Physical-Layer Security Analysis (DEFERRED)}

\textbf{Requires:} Tiers~1 and 2 published.

\textbf{Goal:} Publish a high-impact paper on a QKD side channel that \textit{only
a physical-layer simulator can reveal}.

\subsection*{Tasks}
\begin{enumerate}[nosep]
\item Laser-noise side-channel (RIN resonance $\rightarrow$ amplitude leakage)
\item Finite-key effects with correlated noise (non-i.i.d. RIN/phase noise)
\item Detector side-channel (dead time, afterpulsing, time-shift attacks)
\item Phase modulator flaw analysis (drive ripple $\rightarrow$ deterministic phase errors)
\end{enumerate}

\textbf{Effort: $\sim$5 months.} One strong result beats a catalogue of re-runs.

% ══════════════════════════════════════════════════════════════════
\section{Tier 4 --- Full QKD System Optimisation (OPTIONAL)}

\textbf{Requires:} Tiers~1--3.

\textbf{Goal:} Decoy-state protocol, SKR optimisation, open-source release.

\textbf{Effort: $\sim$10 months.} Strongly diminishing returns after Tier~3.

% ══════════════════════════════════════════════════════════════════
\section{Summary}

\begin{tabularx}{\textwidth}{c L{3cm} L{2cm} L{3cm} L{2.5cm}}
\toprule
Tier & What & Effort & Paper target & Dependency \\
\midrule
0 & Infrastructure & DONE & --- & --- \\
\textbf{1} & \textbf{Software validation} & \textbf{3--4 w} & \textbf{OE\,/\,JLT\,/\,SoftwareX} & \textbf{---} \\
2 & QKD environmental & 6 w & OE\,/\,JLT\,/\,EPJ~QT & Tier~1 \\
3 & QKD security & 5 m & IEEE~T-IFS\,/\,NPJ~QI\,/\,PRX~Q. & Tiers~1--2 \\
4 & Full QKD system & 10 m & Nature Comms\,/\,PRX~Q. & Tiers~1--3 \\
\bottomrule
\end{tabularx}

\subsection*{Recommended Route}
\begin{enumerate}[nosep]
\item \textbf{Tier~1 now} (3--4 weeks). Publish a validated physical-layer
  simulator paper. This is the immediate priority.
\item \textbf{If Tier~2's QBER signal is strong} ($\ge 0.5\%$ delta), publish
  an environmental journal paper (6 weeks).
\item \textbf{Tier~3} (5 months). The laser-noise side-channel is the
  highest-impact result this platform enables.
\item \textbf{Stop} after Tier~3 unless you want a flagship capstone.
  Diminishing returns past this point are steep.
\end{enumerate}

\subsection*{Competitive Landscape}

{\small
\begin{tabularx}{\textwidth}{L{3.2cm} L{1.6cm} L{1.6cm} L{1.4cm} L{1.6cm} L{1.4cm}}
\toprule
Advantage & opto-sim & NetSquid & qkdX & Comsis & Caputo \\
\midrule
Physical-layer field & \done & --- & --- & \done & \done \\
Multi-experiment validated & \textsl{in prog.} & --- & --- & --- & --- \\
Open-source & \done & \done & --- & --- & --- \\
Modern (Python + NumPy) & \done & \done & --- & --- & \done \\
Laser RIN + phase noise & \done & --- & --- & --- & --- \\
Temp/bend birefringence & \done & --- & --- & --- & --- \\
FFT-based CD + PMD & \done & --- & --- & --- & --- \\
APD with excess noise & \done & --- & --- & --- & --- \\
Seeded reproducibility & \done & --- & --- & --- & --- \\
\bottomrule
\end{tabularx}}

\vspace{0.3cm}
\noindent This roadmap positions opto-sim as the first\textbf{ open-source,
multi-experiment-validated, physical-layer fiber-optic simulator} --- a unique
position in the landscape.

\end{document}
